% 397_Bell_Synchronisation_De_ch.tex
% Johann Pascher, 18. Nov. 2026
% Bell-Verletzung durch Synchronisation: Pilotwellen-Tröpfchen und die deterministische Lesart der FFGFT

\providecommand{\BM}{\textbf{[B]}}
\providecommand{\KM}{\textbf{[K]}}
\providecommand{\SM}{\textbf{[S]}}
\providecommand{\XM}{\textbf{[X]}}
\providecommand{\QM}{\textbf{[Q]}}
\providecommand{\EM}{\textbf{[E]}}

\chapter*{Bell-Verletzung durch Synchronisation}
\addcontentsline{toc}{chapter}{Bell-Verletzung durch Synchronisation}

\begin{abstract}
Zwei Arbeiten aus der Pilotwellen-Hydrodynamik zeigen, dass zwei klassische Tröpfchen, die über ein
gemeinsames Wellenfeld gekoppelt sind, die Bell-Ungleichung in der CHSH-Form verletzen können: in
einem statischen Test mit $S=2{,}49$ und in einem reduzierten Modell mit $S\approx3{,}15$ statisch und
bis $S\approx2{,}85$ in einem Test, bei dem die Kopplung vor der Messung abgeschaltet wird. Die
Autoren führen die Verletzung auf deterministische, wellenvermittelte Synchronisation und auf die
Abhängigkeit des Wellenfelds von den Messeinstellungen zurück. Das Dokument stellt die Arbeiten mit
wörtlichen Zitaten dar, rechnet die Grenzen nach und ordnet sie in die deterministische Lesart der
FFGFT ein (Dok.~230, 328, 334). Das Ergebnis hat zwei Seiten. Die Tröpfchen zeigen, dass ein
deterministischer Einzelablauf mit kontextueller Festlegung der Ausgänge, wie ihn Dok.~230 annimmt,
Bell-Verletzungen tragen kann. Sie zeigen aber auch, dass Messabhängigkeit allein keine
Tsirelson-Grenze kennt: Beide Werte des reduzierten Modells liegen über $2\sqrt2$, der statische um
mehr als 200 Standardfehler. Die quantenmechanische Grenze folgt in der FFGFT aus den geometrischen
Quotienten von Dok.~230, nicht aus Synchronisation; dynamische Synchronisation scheidet als
Mechanismus der FFGFT-Korrelation aus, weil die Kopplung $K_{\mathrm{eff}}=2\pi\xi$ rund 6000-fach
schwächer ist als im Tröpfchenmodell.
\end{abstract}

% ============================================================
\section*{Statusmarkierungen}
% ============================================================
\begin{center}
\begin{tabular}{@{}lp{11cm}@{}}
\textbf{[B]} & Algebraisch bewiesen. \\[2pt]
\textbf{[K]} & Numerisch bestätigt (Prüfskript, Vergleich mit Messwerten oder Literaturwerten). \\[2pt]
\textbf{[S]} & Setzung, offen oder Einordnung ohne eigene Rechnung. \\[2pt]
\textbf{[X]} & Ausgeschlossen. \\[2pt]
\textbf{[Q]} & Aus der Literatur übernommen, mit Quelle. \\[2pt]
\textbf{[E]} & Etablierte Physik, im Dokument nicht hergeleitet. \\
\end{tabular}
\end{center}
Prüfskript: \texttt{2/python/Dok397\_Skripte/pruef\_397\_bell\_synchronisation.py}.
Wörtliche Zitate stehen in der Originalsprache in Anführungszeichen und sind mit \QM{} markiert.

% ============================================================
\section{Anlass und Frage}
% ============================================================

Die FFGFT liest die Quantenmechanik deterministisch. Jede Messung ist ein Einzelvorgang, dessen
Ausgang durch die Lage des Systems und eine in der Wechselwirkung mit der Messanordnung festgelegte
Größe $\lambda$ bestimmt ist; der Born-Wert $|\alpha|^2$ ist ein geometrischer Quotient und von Grund
auf keine Wahrscheinlichkeit (Dok.~230). Bell-Korrelationen gelten als real, aber als topologische
Verbindung auf $T^4$, nicht als Fernwirkung im offenen Raum (Dok.~230, 193). Damit steht die FFGFT
in der Reihe der Ansätze, die eine der Voraussetzungen des Bell-Theorems nicht teilen: Die
Festlegung des Ausgangs in der Wechselwirkung mit der Messanordnung macht das Modell kontextuell.

Im August 2026 ist eine Arbeit erschienen, die genau diesen Weg in einem vollständig klassischen
System durchrechnet \QM~\cite{Lopez2026}. Zwei Tröpfchen, die auf einem schwingenden Ölbad laufen und
über ihr gemeinsames Wellenfeld gekoppelt sind, verletzen dort die CHSH-Ungleichung, auch nachdem
die Kopplung vor der Messung abgeschaltet wurde. Die Frage dieses Dokuments ist, was das für die
deterministische Lesart der FFGFT bedeutet: Was stützt es, was nicht, und was lässt sich daraus
für die FFGFT lernen.

% ============================================================
\section{Das System: Tröpfchen, die ihre eigene Welle reiten}
% ============================================================

Couder und Fort haben 2005 gefunden, dass ein millimetergroßes Öltröpfchen auf einem senkrecht
schwingenden Bad aus demselben Öl nicht verschmilzt, sondern springt und dabei bei jedem Aufprall
eine stehende Faraday-Welle anregt \QM~\cite{Couder2005}. Oberhalb einer Schwelle beginnt es
seitlich zu laufen, angetrieben von der Steigung seiner eigenen Welle. Das Tröpfchen ist ein
Teilchen, die Welle ein reales Feld, und beide sind über die Aufprallfolge gekoppelt. Die
Arbeitsgruppe von Bush hat daraus ein eigenes Gebiet gemacht, die Pilotwellen-Hydrodynamik, mit
einer wachsenden Liste klassischer Gegenstücke zu Quanteneffekten: Beugung am Einzel- und Doppelspalt,
quantisierte Bahnen, Tunneln, Statistiken in begrenzten Becken, die den Eigenmoden des Beckens folgen
\QM~\cite{BushOza2021}.

Der entscheidende Zug des Systems ist sein Gedächtnis. Die Welle klingt langsam ab, und an jeder
Stelle trägt sie die Spur aller früheren Aufpralle; die Kraft auf das Tröpfchen hängt deshalb von
seiner ganzen Vorgeschichte ab. Die quantenähnlichen Effekte treten im Bereich langen Gedächtnisses
auf, wenn die Wellen am beständigsten sind \QM~\cite{Lopez2026}. Das System ist eine makroskopische
Verwirklichung der Doppellösung von de Broglie: eine reale Führungswelle und daneben eine aus
vielen Läufen entstehende statistische Welle.

% ============================================================
\section{Zwei Bell-Tests mit Tröpfchen}
% ============================================================

\subsection{Der statische Test von Papatryfonos et al.}

Papatryfonos, Vervoort, Nachbin, Labousse und Bush haben zwei Tröpfchen numerisch simuliert, jedes in
einem Paar von Mulden, zwischen denen es tunneln kann; die beiden Teilsysteme sind über einen
Kopplungskanal verbunden, den die Tröpfchen nicht betreten, durch den aber Wellen laufen
\QM~\cite{Papatryfonos2023}. Der zweiwertige Ausgang ist die Mulde, in der das Tröpfchen zum
Messzeitpunkt liegt (außen $+1$, innen $-1$); die Messeinstellung ist die Tiefe der Barriere
zwischen den Mulden. Den Schlüsselschritt nennen die Autoren selbst: \enquote{the system geometry may
serve as a proxy for analyzer settings} \QM~\cite{Papatryfonos2023}. Mit symmetrischen Einstellungen
fanden sie als größten Wert $S_{\max}=2{,}49\pm0{,}04$, mit einer statistischen Sicherheit über 20
Standardabweichungen, allerdings nur in einem schmalen Bereich der Barrierentiefen \QM. Die
Verletzung führen sie auf die wellenvermittelte Kopplung zurück: Die Form der Pilotwelle und damit
die Lage der Tröpfchen hängt von der Geometrie beider Teilsysteme ab. Ein dynamischer Test, bei dem
die Kopplung vor dem Umstellen der Einstellungen abgeschaltet wird, war damals noch offen.

\subsection{Das reduzierte Modell von López, Valani, Li und Bush}

Die Arbeit vom August 2026 ersetzt die volle Wellengleichung durch einen sinusförmigen Wellenkern
$W(x)=\cos x$, der eine exakte Rückführung auf gewöhnliche Differentialgleichungen erlaubt
\QM~\cite{Lopez2026}. Jedes Tröpfchen sitzt in einer Doppelmulde; seine Bewegung folgt aus Trägheit,
Reibung, der Kraft der eigenen vergangenen Welle, der Kraft der Welle des Partners und der äußeren
Mulde. Mit Gedächtnisvariablen $Y_{ij}$, $Z_{ij}$, die die exponentiell gewichtete Vorgeschichte
tragen, entsteht ein zwölfdimensionales System aus zwei gekoppelten Lorenz-ähnlichen Teilsystemen.
Ein Schalter $\zeta\in\{0,1\}$ schaltet die Kopplung zwischen den Tröpfchen ein oder aus.

Die Symmetrie des Systems liefert zwei invariante Synchronisationsmannigfaltigkeiten, eine für
gleichlaufende und eine für gegenläufige Bewegung. Auf ihnen laufen die Tröpfchen periodisch oder
chaotisch synchron; im chaotischen Bereich wechselt die Bahn unregelmäßig zwischen beiden
Mannigfaltigkeiten. Je nach Mulde, also je nach Einstellung, hält sich die Bahn länger in der einen
oder der anderen auf, und daraus entstehen die Korrelationen. Im statischen Test konvergiert der
Mittelwert über 1000 zufällige Anfangsbedingungen auf $\langle S\rangle\approx3{,}15$ bei einem
Standardfehler von etwa $1{,}5\times10^{-3}$ \QM~\cite{Lopez2026}. Die Verletzungen liegen nicht in
einer Nische, sondern in einem breiten zusammenhängenden Bereich mittleren Gedächtnisses und
mittlerer Wellenamplitude; am stärksten sind sie im periodischen Bereich.

\subsection{Der dynamische Test}

Der neue Schritt ist ein Test in drei Stufen \QM~\cite{Lopez2026}. Zuerst laufen die Tröpfchen mit
eingeschalteter Kopplung und den Einstellungen $a$, $a^*$, bis sie synchron sind. Dann wird die
Kopplung abgeschaltet ($\zeta=0$), und beide Teilsysteme laufen eine Zeit $T_2$ unabhängig weiter.
Schließlich werden die Einstellungen auf $b$, $b^*$ umgestellt, indem die Barrieren erhöht werden;
das treibt jedes Tröpfchen schnell in eine der beiden Mulden, und diese Mulde ist der Ausgang. Der
Mittelwert $\langle|S|\rangle$ schwingt mit $T_2$ und überschreitet die Grenze 2 periodisch; das
Maximum liegt bei etwa $2{,}85$ \QM. Die Autoren fassen das Ergebnis so zusammen, dass
\enquote{nonlinear deterministic dynamics can produce Bell violations through wave-mediated
synchronization mechanisms} \QM~\cite{Lopez2026}, und an anderer Stelle kürzer: \enquote{memory
(specifically, time-delay coupling) can account for entanglement} \QM~\cite{Lopez2026}.

Wichtig ist, wie die Autoren die Verletzung selbst begründen. Die verborgenen Größen sind nicht nur
Ort und Geschwindigkeit der Tröpfchen, sondern auch das ausgedehnte Wellenfeld, und dessen Form
hängt von der Geometrie der Mulden ab. Die Verletzungen seien deshalb so zu verstehen, dass
\enquote{the wave form is influenced by the measurement settings} \QM~\cite{Lopez2026}. Die
Voraussetzung der Messunabhängigkeit, $p(\lambda\,|\,a,b)=p(\lambda)$, ist im Tröpfchensystem
also nicht erfüllt, und die Autoren sagen das ausdrücklich. Vervoort hatte schon früher gefragt,
ob diese Voraussetzung in Systemen mit einem Hintergrundfeld überhaupt gelten kann
\QM~\cite{Vervoort2018}.

% ============================================================
\section{Nachrechnung der Grenzen}
% ============================================================

Für die Einordnung genügen drei Grenzen des CHSH-Werts $S=E(a,b)+E(a',b)+E(a,b')-E(a',b')$:
die Bell-Grenze $|S|\le2$ für lokale Modelle mit Messunabhängigkeit \EM~\cite{Bell1964,CHSH1969},
die Tsirelson-Grenze $|S|\le2\sqrt2=2{,}828$ für Quantenkorrelationen \EM~\cite{Tsirelson1980} und
die algebraische Grenze $|S|\le4$, die aus $|E|\le1$ folgt.

\begin{center}
\small
\begin{tabular}{@{}p{6.2cm}p{1.4cm}p{6.6cm}@{}}
\toprule
\textbf{Wert} & \textbf{$S$} & \textbf{Lage} \\
\midrule
Papatryfonos et al., statisch & $2{,}49$ & über 2, unter $2\sqrt2$ \KM \\
López et al., dynamisch (Maximum) & $2{,}85$ & $0{,}022$ über $2\sqrt2$, also $0{,}76\,\%$ \KM \\
López et al., statisch (Mittel) & $3{,}15$ & $0{,}32$ über $2\sqrt2$, rund 214 Standardfehler \KM \\
\bottomrule
\end{tabular}
\end{center}

Die Werte des reduzierten Modells liegen über der Tsirelson-Grenze. Die Autoren vermerken das beim
statischen Wert selbst. Daraus folgt eine Aussage, die für jede deterministische Lesart wichtig ist:
Die Tröpfchenkorrelationen sind keine Quantenkorrelationen, sondern stärkere. Das ist kein Mangel
der Arbeit, sondern liegt am Mechanismus. Ist die Verteilung der verborgenen Größen von den
Einstellungen abhängig, so lässt sich jede Korrelationstabelle erzeugen. Das Prüfskript zeigt das an
einem deterministischen Modell, das für das Paar $(a',b')$ ungleiche und sonst gleiche Ausgänge
liefert; es erreicht die algebraische Grenze $S=4$ \KM. Messabhängigkeit erklärt also, warum die
Bell-Grenze 2 fallen kann. Sie erklärt nicht, warum die Natur bei $2\sqrt2$ stehen bleibt.

Zum dynamischen Test eine Anmerkung zur Reichweite. Während der Kopplungsphase stehen die
Einstellungen $a$, $a^*$ bereits fest, und sie formen die Welle und den synchronen Zustand, der nach
dem Abschalten weiterläuft. Die späteren Einstellungen $b$, $b^*$ sind mit $a$, $a^*$ fest gepaart.
Der Test schließt damit die Signalübertragung nach dem Abschalten aus, nicht aber die Abhängigkeit
der verborgenen Größen von den Einstellungen; das entspricht der Lesart der Autoren. Er ist deshalb
kein Gegenstück zu den lückenlosen Bell-Tests mit zufällig und spät gewählten Einstellungen
\EM~\cite{Hensen2015}, sondern ein Beispiel dafür, wie Messabhängigkeit in einem klassischen Feld
konkret aussieht \SM.

% ============================================================
\section{Was die Arbeiten zeigen}
% ============================================================

Erstens: Ein vollständig deterministisches, klassisches System mit realem Feld kann die
Bell-Ungleichung verletzen, ohne Fernwirkung und ohne Zufall im Einzelablauf. Der Zufall der
Messreihe stammt aus dem Chaos der Bahn und aus den Anfangsbedingungen, so wie beim Würfel; die
Autoren ziehen diesen Vergleich selbst \QM~\cite{Lopez2026}. Das ist eine Existenzaussage, und für sie
genügt ein Modell.

Zweitens: Der Preis ist die Messabhängigkeit. Ohne sie gilt die Bell-Grenze für jedes lokale Modell,
und kein Mechanismus kann sie überschreiten \EM. Mit ihr ist die Grenze offen bis 4. Diesen Ausweg
nennt man Superdeterminismus \QM~\cite{HossenfelderPalmer2020}. Palmer, der
den Superdeterminismus seit Jahren vertritt, schreibt dazu nüchtern: \enquote{No experiment to date
has closed the superdeterministic \enquote{loophole}} \QM~\cite{Palmer2023}. Die Tröpfchen schließen diese
Lücke nicht, sie geben ihr eine anschauliche Mechanik.

Drittens: Die Stärke der Korrelation ist im Tröpfchensystem eine Frage der Parameter, nicht eines
Prinzips. Das reduzierte Modell überschreitet Tsirelson; das volle Modell von Papatryfonos et al.
bleibt darunter. Eine Grenze bei $2\sqrt2$ entsteht in keinem der beiden Modelle aus der Struktur.

% ============================================================
\section{Einordnung in die FFGFT}
% ============================================================

\subsection{Gemeinsamkeiten}

Die Messung in der FFGFT ist ein deterministischer Polübergang (Dok.~230): Der Ausgang ist
$A(z,\lambda)=\operatorname{sgn}(z-\lambda)$, festgelegt durch die Lage $z$ des Systems und die
Größe $\lambda$ der Messanordnung, und diese Festlegung geschieht in der Wechselwirkung zwischen
System und Messanordnung. Das ist dieselbe logische Stelle, an der die Tröpfchen die
Messunabhängigkeit verlassen: Die Einstellung formt die Größe, die den Ausgang bestimmt. Dok.~230
nennt das Modell deshalb kontextuell, und das Kochen--Specker-Theorem greift nicht. Die Tröpfchen
liefern damit ein konkretes klassisches Beispiel für die Art von Kontextualität, die die FFGFT
annimmt \SM.

Auch die Lesart des Zufalls stimmt überein. In beiden Fällen ist der Einzelablauf festgelegt, und
die Häufigkeiten der Messreihe entstehen, weil die bestimmende Größe nicht kontrolliert wird. In
Dok.~230 betrifft die Unkenntnis nur $\lambda$; bei den Tröpfchen betrifft sie die Anfangsbedingungen
und das Wellenfeld. Und schließlich ist Messung in beiden Fällen Kopplung: Dok.~328 beschreibt jede
Messung als Ankopplung eines Auslesesystems mit einem Kopplungsgrad $\kappa$, und die Tröpfchen
messen, indem eine Barriere erhöht wird, bis das System in eine Mulde fällt.

\subsection{Unterschiede}

\textbf{Kopplungsstärke.} Im Tröpfchenmodell ist die dimensionslose Wellenamplitude von der
Größenordnung 5 bei einer Dämpfung von 1; die Kopplung ist stark, und Synchronisation ist das
Ergebnis. In der FFGFT ist die effektive Kopplung in der Kreisabbildung $K_{\mathrm{eff}}=2\pi\xi
\approx8{,}4\times10^{-4}$, die Arnold-Zunge der Ordnung $1/2$ ist $5{,}6\times10^{-8}$ breit
(Dok.~328) \KM. Das Verhältnis beträgt rund 6000 \KM. Dynamische Synchronisation, wie sie die
Tröpfchen tragen, ist in der FFGFT bei dieser Kopplung als Mechanismus der Korrelation
ausgeschlossen \XM. Das deckt sich mit Dok.~328 und 343: Die Diskretheit der FFGFT ist topologisch,
nicht dynamisch. Die Bell-Verbindung ist in der FFGFT eine topologisch verbundene
Zylinderkonfiguration auf $T^4$ (Dok.~230), kein eingerasteter Zustand zweier Oszillatoren.

\textbf{Die Tsirelson-Grenze.} Hier liegt der wichtigste Unterschied. Die Quotienten aus Dok.~230
liefern für das Singulett den Anteil gleicher Ausgänge $\sin^2\!\big((a-b)/2\big)$ und daraus
$E(a,b)=-\cos(a-b)$ \KM. Diese Korrelation erreicht bei den Winkeln $0$, $90^\circ$, $45^\circ$,
$-45^\circ$ genau $2\sqrt2$ und überschreitet sie bei keiner Wahl der Winkel; eine Gittersuche im
Prüfskript bestätigt das \KM. Die Grenze ist hier keine zusätzliche Annahme, sondern das Maximum
einer Kosinuskombination. Wo die Tröpfchen über $2\sqrt2$ hinausgehen, weil Messabhängigkeit allein
nichts begrenzt, bleibt die FFGFT bei der Grenze, weil ihre Gewichte geometrische Quotienten sind
(Dok.~230, 367, 370).

\textbf{Was dafür offen ist.} Das lokale Modell $A=\operatorname{sgn}(z-\lambda)$ mit einem
gemeinsamen, von den Einstellungen unabhängigen $\lambda$ ergibt für zwei Teilsysteme nur die
Sägezahnkorrelation $E=-1+2|a-b|/\pi$ und damit höchstens $S=2$ \KM. Die Kosinusform verlangt mehr:
die topologische Verbindung auf $T^4$, die die beiden Zylinderkonfigurationen nicht unabhängig sein
lässt. Im Korpus wird die Übereinstimmung mit $-\cos(a-b)$ über die vollständige Übersetzbarkeit
zwischen Hilbertraum und Zylinderdarstellung getragen (Dok.~230); eine Herleitung der Kosinusform
aus der topologischen Verbindung selbst, ohne Rückgriff auf den Hilbertraum, liegt nicht vor \SM.
Das ist die Stelle, an der die Tröpfchen der FFGFT eine genaue Aufgabe stellen: zu zeigen, dass die
Verbindung auf $T^4$ genau die Kosinusform erzeugt und nicht mehr.

\textbf{Gedächtnis und Zeit.} Die Tröpfchen brauchen ein Gedächtnis von endlicher Dauer, und ihre
Korrelation schwingt mit der Wartezeit $T_2$. In der FFGFT trägt jede Mode ihre eigene Zeitskala
$\tilde T=1/m$ (Dok.~334); die Korrelation ist keine gespeicherte Vorgeschichte, sondern eine
Eigenschaft der Konfiguration. Ein Schwingen der Bell-Korrelation mit der Wartezeit zwischen
Entstehung und Messung erwartet die FFGFT daher nicht \SM.

\subsection{Was die FFGFT daraus mitnimmt}

Die Tröpfchen schärfen die Frage, an der sich eine deterministische Lesart messen lassen muss. Dass
Bell-Verletzungen ohne Fernwirkung und ohne Zufall möglich sind, ist mit ihnen anschaulich gezeigt.
Die eigentliche Frage ist, warum die Natur genau bei $2\sqrt2$ aufhört. Eine Lesart, die nur
Messabhängigkeit anbietet, beantwortet sie nicht. Die FFGFT hat dafür eine Antwort im Ansatz: die
Born-Gewichte als geometrische Quotienten einer geschlossenen Geometrie. Ausgearbeitet ist die
Antwort für den Einzelfall (Dok.~230, 367); für das verschränkte Paar steht sie aus \SM.

% ============================================================
\section{Was nicht übereinstimmt}
% ============================================================

\begin{itemize}[leftmargin=2em,itemsep=3pt,topsep=3pt]
  \item Die Tröpfchen leben in einer Dimension und in einem Modell mit idealisiertem Wellenkern;
        die Autoren nennen selbst als offen, ob die Verletzungen in einem vollständig dynamischen,
        lokalen Wellenmodell mit endlicher Laufzeit bestehen bleiben \QM~\cite{Lopez2026}.
  \item Die Ergebnisse sind Simulationen; ein Experiment mit echten Tröpfchen gibt es bisher nicht.
  \item Die Messeinstellung der Tröpfchen ist eine Geometrie, die Welle und Bahn formt; in der FFGFT
        ist $\lambda$ eine Größe der Messanordnung auf derselben Achse wie $z$, ohne ausgedehntes Feld
        im Ortsraum.
  \item Die Tröpfchen überschreiten Tsirelson, die FFGFT nicht.
  \item Die Tröpfchen synchronisieren dynamisch bei starker Kopplung; die FFGFT ist bei
        $K_{\mathrm{eff}}=2\pi\xi$ tief unterkritisch.
\end{itemize}

% ============================================================
\section{Übersicht}
% ============================================================
\begin{center}
\small
\begin{tabular}{@{}p{8.2cm}p{1.2cm}p{4.4cm}@{}}
\toprule
\textbf{Aussage} & \textbf{Status} & \textbf{Grund} \\
\midrule
Tröpfchen verletzen CHSH: $2{,}49$; $3{,}15$; dynamisch bis $2{,}85$ & \QM & \cite{Papatryfonos2023,Lopez2026} \\
Begründung über Messabhängigkeit der Welle & \QM & \cite{Lopez2026} \\
$3{,}15$ und $2{,}85$ liegen über $2\sqrt2$ & \KM & Prüfskript \\
Messabhängigkeit erlaubt $S$ bis 4 & \KM & Prüfskript \\
Born-Quotienten aus Dok.~230 ergeben $-\cos(a-b)$, Maximum genau $2\sqrt2$ & \KM & Prüfskript \\
$\operatorname{sgn}(z-\lambda)$ mit gemeinsamem $\lambda$: Sägezahn, $S\le2$ & \KM & Prüfskript \\
Kontextuelle Festlegung in Dok.~230 entspricht der Messabhängigkeit der Tröpfchen & \SM & Einordnung \\
Synchronisation als FFGFT-Mechanismus bei $K_{\mathrm{eff}}=2\pi\xi$ & \XM & Dok.~328, Prüfskript \\
Kosinusform aus der topologischen Verbindung auf $T^4$ hergeleitet & \SM & offen \\
kein Schwingen der Korrelation mit der Wartezeit & \SM & Dok.~334 \\
\bottomrule
\end{tabular}
\end{center}

\begingroup\renewcommand{\chapter}[2]{}\small
\begin{thebibliography}{99}
\bibitem{Lopez2026} Á.~G.~López, R.~N.~Valani, Y.~Li, J.~W.~M.~Bush, \textit{Synchronization induces Bell violations in a model of walking droplets}, arXiv:2608.21915 (2026).
\bibitem{Papatryfonos2023} K.~Papatryfonos, L.~Vervoort, A.~Nachbin, M.~Labousse, J.~W.~M.~Bush, \textit{A platform for investigating Bell correlations in pilot-wave hydrodynamics}, arXiv:2208.08940v2 (2023).
\bibitem{Couder2005} Y.~Couder, S.~Protière, E.~Fort, A.~Boudaoud, \textit{Walking and orbiting droplets}, Nature \textbf{437} (2005) 208.
\bibitem{BushOza2021} J.~W.~M.~Bush, A.~U.~Oza, \textit{Hydrodynamic quantum analogs}, Reports on Progress in Physics \textbf{84} (2021) 017001.
\bibitem{Vervoort2018} L.~Vervoort, \textit{Are hidden-variable theories for pilot-wave systems possible?}, Foundations of Physics \textbf{48} (2018) 803.
\bibitem{Bell1964} J.~S.~Bell, \textit{On the Einstein Podolsky Rosen paradox}, Physics \textbf{1} (1964) 195.
\bibitem{CHSH1969} J.~F.~Clauser, M.~A.~Horne, A.~Shimony, R.~A.~Holt, \textit{Proposed experiment to test local hidden-variable theories}, Physical Review Letters \textbf{23} (1969) 880.
\bibitem{Tsirelson1980} B.~S.~Cirel'son, \textit{Quantum generalizations of Bell's inequality}, Letters in Mathematical Physics \textbf{4} (1980) 93.
\bibitem{Hensen2015} B.~Hensen et al., \textit{Loophole-free Bell inequality violation using electron spins separated by 1.3 kilometres}, Nature \textbf{526} (2015) 682.
\bibitem{HossenfelderPalmer2020} S.~Hossenfelder, T.~N.~Palmer, \textit{Rethinking superdeterminism}, Frontiers in Physics \textbf{8} (2020) 139.
\bibitem{Palmer2023} T.~N.~Palmer, \textit{Superdeterministic loophole}, Leserbrief, Physics Today, Juni 2023.
\end{thebibliography}
\endgroup
